\usepackage{xeCJK}
\usepackage{geometry}
\geometry{left=2cm,right=2cm,top=2.5cm,bottom=2.5cm}
\setlength{\headheight}{12.7pt}

\usepackage{titlesec}
\titleformat{\section}
  {\normalfont\large\bfseries}
  {\thesection}
  {1em}
  {}
\titlespacing*{\section}
  {0pt}
  {3.5ex plus 1ex minus .2ex}
  {2.3ex plus .2ex}

\usepackage{amsmath}
\usepackage{amssymb}
\usepackage{amsthm}
\usepackage{enumitem}
\setlist[enumerate]{itemsep=0pt,topsep=0pt,parsep=0pt,leftmargin=0pt,itemindent=2.5em}
\setlist[itemize]{itemsep=0pt,topsep=0pt,parsep=0pt,leftmargin=0pt,itemindent=2.5em}
\usepackage{fancyhdr}
\pagestyle{fancy}
\fancyhf{}
\usepackage{graphicx}
\usepackage{hyperref}
\usepackage{indentfirst}
\usepackage{lmodern}


\begin{document}
\setlength{\abovedisplayskip}{1pt}
\setlength{\belowdisplayskip}{1pt}

\section{字典序}

字典序是多个序关系的结合.

\vspace{10pt}

\hypertarget{sec:定义1.1}{}
\noindent\textbf{定义1.1 (字典序)}

给定\href{04 良序关系#nameddest=sec:定义1.1}{良序结构} $(I,\leqslant)$和
一族集合$(f_i)_{i\in I}$, 再给定每个$f_i$上的序$\leqslant_i$.
定义$\displaystyle\prod_{i\in I}f_i$上的\textbf{字典序}:
\[\begin{aligned}
    &p<_{\mathrm{lex}}q\quad\overset{\mathrm{def}}{\leftrightarrow}\quad&&
    \exists i\in I:\,\,(\,\,\forall k<i:p(k)=q(k)\,\,\land\,\,p(i)<_iq(i)\,\,).
    \\[-3pt]
    &p\leqslant_{\mathrm{lex}}q\quad\overset{\mathrm{def}}{\leftrightarrow}\quad
    &&p<_{\mathrm{lex}}q\,\,\lor\,\,p=q.
\end{aligned}\]

\vspace{10pt}

字典序会保持原有的偏序和全序性质.

\vspace{10pt}

\hypertarget{sec:定理1.1}{}
\noindent\textbf{定理1.1}

给定良序结构$(I,\leqslant)$和一族集合$(f_i)_{i\in I}$,
再给定每个$f_i$上的序$\leqslant_i$, 记字典序为$\leqslant_{\mathrm{lex}}$.

\begin{enumerate}
    \item 如果每个$\leqslant_i$都是偏序, 那么$\leqslant_{\mathrm{lex}}$也是偏序,
    \item 如果每个$\leqslant_i$都是全序, 那么$\leqslant_{\mathrm{lex}}$也是全序.
\end{enumerate}

\noindent{\kaishu 证明.}

\begin{enumerate}
    \item 反自反性: 显然.
    
    \hspace{2.17em}
    传递性: 如果$p<_{\mathrm{lex}}q$且$q<_{\mathrm{lex}}r$, 那么存在$i,j\in I$使得
    \vspace{3pt}
    \[\left\{\begin{aligned}
        &\forall k<i:p(k)=q(k)&&\quad\land\quad p(i)<_iq(i),\\[-3pt]
        &\forall k<j:q(k)=r(k)&&\quad\land\quad q(j)<_jr(j),
    \end{aligned}\right.\]
    
    \vspace{3pt}\hspace{2.17em}
    如果$i<j$, 那么对任意的$k<i$有$p(k)=q(k)=r(k)$,
    并且$p(i)<_iq(i)=r(i)$, 从而$p<_{\mathrm{lex}}r$ ;

    \hspace{2.17em}
    如果$i=j$, 那么对任意的$k<i$有$p(k)=q(k)=r(k)$,
    并且$p(i)<_iq(i)<_ir(i)$, 从而$p<_{\mathrm{lex}}r$ ;

    \hspace{2.17em}
    如果$i>j$, 那么对任意的$k<j$有$p(k)=q(k)=r(k)$,
    并且$p(j)=q(j)<_jr(j)$, 从而$p<_{\mathrm{lex}}r$ .

    \item 三歧性: 对任意的$p,q$, 如果$p\neq q$, 那么
    \[
        \{i\in I\mid p(i)\neq q(i)\}
    \]
    是$I$的非空子集, 记最小元为$i$, 则对任意的$k<i$有$p(k)=q(k)$.

    \hspace{2.17em}
    此时如果$p(i)<_iq(i)$则$p<_{\mathrm{lex}}q$,
    如果$p(i)>_iq(i)$则$p>_{\mathrm{lex}}q$. \hfill $\qedsymbol$
\end{enumerate}

\vspace{10pt}

但是字典序未必一定保持良序的性质.

举个例子, 考虑$\displaystyle\prod_{i\in\omega}2=\,\!^\omega2$上的字典序,
对任意自然数$n$定义
\vspace{-3pt}
\[
    p_n:\omega\to 2,\quad i\mapsto\left\{\begin{aligned}
        &0,\quad&&i<n,\\[-3pt]
        &1,\quad&&i\geqslant n.
    \end{aligned}\right.
\]
接下来, 对任意自然数$n$有
\[
    \forall i<n:p_{n+1}(i)=0=p_n(i)\quad\land\quad p_{n+1}(n)=0<1=p_n(n),\\[3pt]
\]
所以$p_{n+1}<_{\mathrm{lex}}p_n$. 所以$(p_n)_{n\in\omega}$是一个严格递降序列,
由\href{04 良序关系#nameddest=sec:定理1.2}{引理}可知字典序不是良序.





\newpage

有限Descartes积上的字典序就可以保持良序性质.

\vspace{10pt}

\hypertarget{sec:定理1.2}{}
\noindent\textbf{定理1.2}

给定自然数$n$和一族集合$(f_i)_{i<n}$,
再给定每个$f_i$上的良序$\leqslant_i$, 那么字典序$\leqslant_{\mathrm{lex}}$也是良序.

\noindent{\kaishu 证明.}

任取$\displaystyle\prod_{i<n}f_i$的非空子集$S$, 下证$S$有最小元.
显然$n$非零, 对任意的$i\leqslant n$递归定义$S_i$ :
\vspace{6pt}
\begin{enumerate}
    \item $S_0=S$,
    \item $S_{i+1}=\{p\in S_i\mid\forall\varphi\in S_i:p(i)\leqslant_i\varphi(i)\}$.
\end{enumerate}

\vspace{5pt}\hspace{-1em}{\kaishu 命题1.} 对任意$i\leqslant n$, $S_i$非空.

\vspace{3pt}
归纳证明. 首先$S_0=S$非空; 其次任取$i<n$, 假设$S_i$非空, 那么
\[
    \{\varphi(i)\mid\varphi\in S_i\}
\]
是$f_i$的非空子集, 记最小元为$m$, 则存在$p\in S_i$使得$p(i)=m$, 从而$p\in S_{i+1}$.

\vspace{5pt}\hspace{-1em}{\kaishu 命题2.}
对任意$i\leqslant n$和$p\in S$,
$p\in S_i$当且仅当$\forall k<i\,\,\forall\varphi\in S_k:p(k)\leqslant_k\varphi(k)$.

\vspace{3pt}
归纳证明. 首先$p\in S$和$\forall k<0\cdots$始终成立; 其次任取$i<n$, 假设
\[
    p\in S_i\quad\leftrightarrow\quad
    \forall k<i\,\,\forall\varphi\in S_k:p(k)\leqslant_k\varphi(k),
\]
那么
\[\begin{aligned}
    p\in S_{i+1}
    \iff&p\in S_i\,\,\land\,\,\forall\varphi\in S_i:p(i)\leqslant_i\varphi(i)\\[-3pt]
    \iff&\forall k<i\,\,\forall\varphi\in S_k:p(k)\leqslant_k\varphi(k)
    \quad\land\quad\forall\varphi\in S_i:p(i)\leqslant_i\varphi(i)\\[-3pt]
    \iff&\forall k<i+1\,\,\,\forall\varphi\in S_k:p(k)\leqslant_k\varphi(k).
\end{aligned}\]

\vspace{5pt}\hspace{-1em}{\kaishu 命题3.} $S_n$的元素是$S$的最小元.

\vspace{3pt}
取$p\in S_n$, 根据命题2可知
对任意$k<n$和$\varphi\in S_k$有$p(k)\leqslant_k\varphi(k)$.

再任取$q\in S$, 假设$q<_{\mathrm{lex}}p$, 即存在$i<n$使得
\[
    \forall k<i:q(k)=p(k)\quad\land\quad q(i)<_ip(i).
\]
那么对任意$k<i$和$\varphi\in S_k$有$q(k)=p(k)\leqslant_k\varphi(k)$,
于是$p,q\in S_i$. 进一步$p(i)\leqslant_iq(i)$, 矛盾. \hfill $\qedsymbol$





\vspace{20pt}

\section{有限运算与字典序}

任给序集$a,b$, 把它们的一般Descartes积上的字典序复制到$a\times b$上,
定义为$a\times b$上的字典序.

对任意的$(x,y),(x',y')\in a\times b$,
$(x,y)\leqslant_{\mathrm{lex}}(x',y')$当且仅当
\[
    x<x'\quad\lor\quad(\,x=x'\quad\land\quad y\leqslant y'\,).
\]

\vspace{10pt}

\hypertarget{sec:定义2.1}{}
\noindent\textbf{定义2.1 (有限不交并)}

任给集合$a,b$, 定义$a\sqcup b\overset{\mathrm{def}}{=}
\{0\}\times a\cup\{1\}\times b$.

如果$a,b$都是序集, 则$a\sqcup b\subset 2\times(a\cup b)$上也有字典序.

\vspace{10pt}

多元的情况依此类推.

\end{document}